
\documentclass{article}
\usepackage{colortbl}
\usepackage{makecell}
\usepackage{multirow}
\usepackage{supertabular}

\begin{document}

\newcounter{utterance}

\centering \large Interaction Transcript for game `cladder', experiment `full\_v1.5\_default', episode 4237 with score{-}Qwen3.5{-}9B{-}sp.
\vspace{24pt}

{ \footnotesize  \setcounter{utterance}{1}
\setlength{\tabcolsep}{0pt}
\begin{supertabular}{c@{$\;$}|p{.15\linewidth}@{}p{.15\linewidth}p{.15\linewidth}p{.15\linewidth}p{.15\linewidth}p{.15\linewidth}}
    \# & \multicolumn{2}{c}{Player} && \multicolumn{2}{c}{Game Master} \\
    \hline

    \theutterance \stepcounter{utterance}  
    & & & \multicolumn{4}{p{0.6\linewidth}}{
        \cellcolor[rgb]{0.9,0.9,0.9}{
            \makecell[{{p{\linewidth}}}]{
                \texttt{\tiny{[P1$\langle$GM]}}
                \texttt{You are an expert in causal inference. The following question is not a typical commonsense query, but rather a meticulously designed question created by a professor specializing in causal inference, intended to assess the students' mastery of the course content.} \\
\texttt{Imagine a self{-}contained, hypothetical world with only the following conditions, and without any unmentioned factors or causal relationships: Kidney stone size has a direct effect on treatment and lip thickness. Treatment has a direct effect on lip thickness. For patients who have small kidney stones and not receiving treatment, the probability of thick lips is 6\%. For patients who have small kidney stones and receiving treatment, the probability of thick lips is 38\%. For patients who have large kidney stones and not receiving treatment, the probability of thick lips is 63\%. For patients who have large kidney stones and receiving treatment, the probability of thick lips is 95\%. The overall probability of large kidney stone is 50\%. For patients receiving treatment, would it be more likely to see thick lips if the patient had received no treatment?} \\
\texttt{Start your answer with 'Yes' or 'No', followed by additional reasoning or evidence to support your explanation.} \\
            }
        }
    }
     \\ \\

    \theutterance \stepcounter{utterance}  
    & \multicolumn{4}{p{0.6\linewidth}}{
        \cellcolor[rgb]{0.561,0.737,0.561}{
            \makecell[{{p{\linewidth}}}]{
                \texttt{\tiny{[P1$\rangle$GM]}}
                \texttt{No} \\
\\ 
\texttt{**Reasoning:**} \\
\\ 
\texttt{To answer this question, we must determine the causal effect of treatment on lip thickness. Specifically, we need to compare the probability of having thick lips under the intervention of receiving treatment (\$P(\textbackslash{}text\{Thick\} | \textbackslash{}text\{Treat\}=1)\$) versus the probability of having thick lips under the intervention of not receiving treatment (\$P(\textbackslash{}text\{Thick\} | \textbackslash{}text\{Treat\}=0)\$).} \\
\\ 
\texttt{The problem provides the conditional probabilities of thick lips given both the kidney stone size (\$S\$) and the treatment status (\$T\$):} \\
\texttt{1.  \$P(\textbackslash{}text\{Thick\} | S=\textbackslash{}text\{small\}, T=0) = 0.06\$} \\
\texttt{2.  \$P(\textbackslash{}text\{Thick\} | S=\textbackslash{}text\{small\}, T=1) = 0.38\$} \\
\texttt{3.  \$P(\textbackslash{}text\{Thick\} | S=\textbackslash{}text\{large\}, T=0) = 0.63\$} \\
\texttt{4.  \$P(\textbackslash{}text\{Thick\} | S=\textbackslash{}text\{large\}, T=1) = 0.95\$} \\
\\ 
\texttt{We are also given the marginal probability of having a large kidney stone:} \\
\texttt{*   \$P(S=\textbackslash{}text\{large\}) = 0.50\$} \\
\texttt{*   Therefore, \$P(S=\textbackslash{}text\{small\}) = 1 {-} 0.50 = 0.50\$} \\
            }
        }
    }
    & & \\ \\

    \theutterance \stepcounter{utterance}  
    & & & \multicolumn{2}{p{0.3\linewidth}}{
        \cellcolor[rgb]{0.9,0.9,0.9}{
            \makecell[{{p{\linewidth}}}]{
                \texttt{\tiny{[GM$|$GM]}}
                \texttt{No} \\
            }
        }
    }
    & & \\ \\

    \theutterance \stepcounter{utterance}  
    & & & \multicolumn{2}{p{0.3\linewidth}}{
        \cellcolor[rgb]{0.9,0.9,0.9}{
            \makecell[{{p{\linewidth}}}]{
                \texttt{\tiny{[GM$|$GM]}}
                \texttt{no} \\
            }
        }
    }
    & & \\ \\

    \theutterance \stepcounter{utterance}  
    & & & \multicolumn{2}{p{0.3\linewidth}}{
        \cellcolor[rgb]{0.9,0.9,0.9}{
            \makecell[{{p{\linewidth}}}]{
                \texttt{\tiny{[GM$|$GM]}}
                \texttt{game\_result = WIN} \\
            }
        }
    }
    & & \\ \\

\end{supertabular}
}

\end{document}
